\section{Conclusion}

\begin{frame}{Summary of Contributions}
  \begin{enumerate}
    \item \textbf{Capacity-aware Observer/Operator architecture} --- structural read/write separation with tool partitioning (29 read-only vs.\ 40 action tools) and mandatory HITL confirmation; +6.1\,pp over monolithic baseline, +10.2\,pp on routing-neutral cases
    \item \textbf{3,135-case Slurm benchmark} --- 11 categories $\times$ 5 scenarios with ground-truth labels for tool use, routing, HITL, and state transitions; publicly released as a reusable artifact
    \item \textbf{Role-scoped fine-tuning methodology} --- QLoRA on Qwen2.5-14B-Instruct distilled from commercial traces with handoff-split, role-filtered training data; \textbf{88.3\% mean weighted score}, closes 29\% of the gap to GPT-5-mini, locally deployable (no API dependency)
  \end{enumerate}
\end{frame}

\begin{frame}{Limitations \& Future Work}
  \begin{columns}[T]
    \begin{column}{0.48\textwidth}
      \textbf{Limitations}
      \begin{itemize}
        \item Synthetic benchmark only — no live-cluster user study
        \item Routing metric (25\%) structurally disadvantages monolithic baseline
        \item Training data from finite mock state space
        \item Single-turn evaluation only
        \item No multi-tenant auth / audit logging
      \end{itemize}
    \end{column}
    \begin{column}{0.48\textwidth}
      \textbf{Future Work}
      \begin{itemize}
        \item Real-cluster validation with live scheduler variability
        \item Multi-turn conversation benchmark
        \item Role-aware authentication \& audit logging
        \item Larger base model (70B) for harder categories
        \item Practitioner user studies
      \end{itemize}
    \end{column}
  \end{columns}
\end{frame}
